\documentclass[12pt]{article}
\usepackage[utf8]{inputenc}
\usepackage[T1]{fontenc}
\usepackage{hyperref}
\title{Sayula Popoluca Language Resources for K-12 Curriculum Development}
\author{Compiled by Claude Sonnet 5.5 (claude-sonnet-5-5) for Andr\'as Kornai}
\date{}
\begin{document}
\maketitle
\section{Introduction}
Sayula Popoluca (ISO 639-3 pos) is a Mixe-Zoquean language of Veracruz, Mexico; estimates range from about 3,000 to 4,800 speakers. Published work is largely by Lawrence and Nancy Clark (SIL/ILV, 1960s) and is dated; two orthographies appear in the literature.
\section{Primary Educational Resources}
\begin{itemize}
\item \emph{Cartilla popoluca de Sayula} (ILV, 1960, 68 pp.) and \emph{Vocabulario popoluca de Sayula} (Clark \& Clark, 1960, 165 pp.) are held by the El Colegio de M\'exico repository; the two catalogue records below cover these titles: \url{https://repositorio.colmex.mx/concern/books/3b591d18n?locale=en} and \url{https://repositorio.colmex.mx/concern/books/jq085m825?locale=en}
\end{itemize}
\section{Linguistic and Reference Materials}
\begin{itemize}
\item Endangered Languages Project bibliography for Sayula Popoluca (includes Clark's \emph{Sayula Popoluca texts, with grammatical outline}, 1961): \url{https://endangeredlanguages.com/lang/5286/bibliography}
\item Glottolog reference, Clark and Clark 1960: \url{https://glottolog.org/resource/reference/id/13662}
\item PHOIBLE source, Clark 1995: \url{https://phoible.org/sources/pos_clark1995}
\item Wikipedia article: \url{https://en.wikipedia.com/wiki/Sayultec}
\end{itemize}
\section{Religious and Community Resources}
None found.
\section{Digital and Audio Resources}
None found.
\section{Curriculum Development Context}
Existing materials are mid-twentieth-century missionary-linguist work; orthography and vocabulary should be checked with current speakers before use in schools.
\end{document}